\scanpage{110}
\mspage{101}
% archive sheet number: 110

plus grand élément dans $L$, soit \deleted{$\unsure{\delta}$} $\Theta$ \deleted{l'ens} \unsure{la} partie de $L$, de cardinal $\{0,1,2\}$, formée desdits éléments. Alors $\unsure{\Phi|\varphi}$ est ouverte sss
\[
\deleted{\alpha=\partial D}\quad\operatorname{card}D\geq2\text{ et }\alpha=\partial\unsure{\mathbb{D}|D}-(\partial D\cap\Theta).
\]

3) Pour \unsure{tout|toute} \deleted{\unsure{suite de préfigures connexes,}}\inserted{\deleted{\unsure{de}} drapeau}
\[
D=\{\Phi_1\overset{\unsure{!|'}}{<}\Phi_2\overset{!}{<}\cdots\overset{!}{<}\Phi_n\}
\]
dans l'ens ordonné des préfigures \unsure{0-connexes},
\[
D\in\operatorname{Drap}_{<^!}^{*}(\operatorname{Préfig}_{\mathrm{conn}}(\unsure{\mathscr{M}|\mathcal{M}})),
\]
soit
\[
\Phi_D\equiv\bigcup_{\unsure{\Psi|\Phi}\in D}\Phi=\bigcup_{1\leq i\leq n}\Phi_i.
\]
C'est une réunion disjointe, et \deleted{la} $\Phi$ est une préfigure, admettant les $\Phi\in D$ comme composantes connexes.

L'application
\[
D\longmapsto\Phi_D=\bigcup_{\Phi\in D}\Phi
\]
\[
\operatorname{Drap}_{<^!}^{*}(\operatorname{Préfig}_{\mathrm{conn}}(\mathscr{M}))\longrightarrow\operatorname{Préfig}(\mathscr{M})
\]
est bijective. \deleted{Pour}

4) Pour que $\Phi_D$ soit fermée, il f. et
